% Section content template for individual Educative lessons
% This file contains the actual content for a specific section
% Generated from Educative API section content

% Section: Quiz on the Sharded Counters' Design
% Section ID: 5659634509021184
% Section Slug: quiz-on-the-sharded-counters-design
% Generated: 2025-12-07T22:17:30.439248



\subsection*{Quiz}
\vspace{0.3cm}
\noindent
\textbf{Question 1:} What kind of problem would occur if our system created a counter with a few shards for a post from a user with one million followers?
\\[0.2cm]
\begin{itemize}
 \item High read contention
 \item \textbf{[CORRECT]} High write contention
 \item Users face delays on a read request
 \item Users get a quick response on the liked post
\end{itemize}
\vspace{0.3cm}
\noindent
\textbf{Question 2:} \textbf{(Select all that apply.)} What problem might arise if shards are selected in an order (sequentially) rather than randomly?
\\
\textit{(Select all that apply.)}
\\[0.2cm]
\begin{itemize}
 \item \textbf{[CORRECT]} The write request queue increases.
 \item \textbf{[CORRECT]} Maximum shard utilization decreases.
 \item One-by-one selection of shards will take more time.
\\[0.1cm]
\textit{Explanation:} 
This statement is not true because only one increment will give the next shard. The time complexity of selecting shards remains O(1).
 \item \textbf{[CORRECT]} The user will not get quick responses.
\end{itemize}
\vspace{0.3cm}
\noindent
\textbf{Question 3:} \textbf{(Select all that apply.)} Which situation needs a larger number of sharded counters to handle the traffic on YouTube?
\\
\textit{(Select all that apply.)}
\\[0.2cm]
\begin{itemize}
 \item \textbf{[CORRECT]} A video posted by a channel with millions of subscribers
 \item A video with a few likes over a long period of time
 \item \textbf{[CORRECT]} A video with millions of views but a few dislikes
 \item Long videos posted by a channel with a few subscribers
\end{itemize}

% End of section content